\chapter*{Introduction}
\label{chap:intro}

\epigraph{\itshape If you look at a multi-player game, it's the people who are playing the game who are often more valuable than all of the animations and models and game logic that's associated with it.}{-- Gabe NEWELL}

\addcontentsline{toc}{chapter}{Introduction}  
\fancyhead[LE,RO]{Introduction}

\section{Context}

Online competitive video games, particularly \glspl{FPS}, now occupy a central place in both the video game industry and the \glspl{esport} ecosystem. Their model relies on technical fairness between players: each participant must be subject to the same rules, mechanical constraints, and informational limitations.

This is nowhere more visible than in \gls{Counter-Strike 2} (\gls{CS2}), \gls{Valve}'s flagship tactical shooter and one of the most-watched titles in professional esports. The Counter-Strike series in particular has accumulated a continuous record of adversarial development between cheats and anti-cheat systems since the first VAC release in 2002, with each new engine generation triggering an immediate re-emergence of cheat software within days of launch \cite{Kuzhii2024}. Major tournaments distribute cash prizes reaching several million dollars, and top-tier organizations field full-time professional rosters whose careers - and livelihoods - depend on the perceived integrity of every match. This economic stake is precisely why most high-level CS2 tournaments are played in LAN (Local Area Network) conditions: bringing players into a single physical venue, on hardware controlled by tournament organizers, drastically reduces the attack surface available to cheaters compared to online play, where each participant's machine remains outside the organizers' control. The contrast between these two environments - tightly supervised LAN finals versus the millions of anonymous online matchmaking games played daily - illustrates the scale of the challenge: competitive integrity cannot realistically be guaranteed everywhere through physical supervision alone.

However, cheating represents a growing threat to this balance. The emergence of automated tools such as \glspl{aimbot} and \glspl{wallhack} alters match integrity, undermines competitive credibility, and directly affects player trust as well as the economic viability of affected titles. In a game where a single well-placed shot can decide a round, and a round can decide a series worth a cash prize, even a small number of undetected cheaters in the matchmaking pool is enough to erode confidence in the entire ranking and reward system - both for the casual player grinding toward a higher rank, and for the aspiring professional whose statistics may be scouted by a team.

To address this issue, many publishers have progressively adopted kernel-level anti-\gls{cheat} systems. These solutions operate with very high privileges, allowing them to monitor memory, processes, and certain hardware interactions within the operating system in order to detect manipulation of the game \gls{client}. While often effective against certain forms of software cheating, they also raise significant concerns regarding \gls{cybersecurity}, privacy protection, and user acceptability. The permanent presence of a third-party privileged component increases the system’s attack surface, introduces a critical point of failure and fuels public debate about the proportionality between technical intrusion and its intended objective.

Beyond pure cybersecurity vulnerabilities, \gls{kernel}-level solutions introduce profound privacy and regulatory dilemmas. Operating at the highest privilege layer (Ring 0) theoretically grants these software components unhindered visibility over the user's entire system, including personal files, background applications, and non-game-related data. This continuous monitoring mechanism challenges traditional notions of informed user consent and clashes directly with modern data protection frameworks, such as the GDPR, which mandate strict principles of data minimization and proportionality. Consequently, this technological choice has sparked a sharp divide within the gaming community: while competitive players often tolerate intrusive measures in the name of competitive integrity, casual users increasingly reject what they perceive as a form of digital surveillance, ultimately threatening game adoption and brand reputation.

\section{Problematic}

In this context, the following question arises: can cheating be effectively detected without compromising system security through kernel-level access? This thesis explores an alternative approach primarily based on \gls{server}-side behavioral analysis and machine learning techniques. Instead of directly inspecting the player’s machine, the idea is to leverage gameplay data itself-match statistics, action sequences, and temporal and mechanical consistency-to identify anomalies characteristic of assisted behavior.

Recent research suggests that analyzing replays and gameplay data in games such as Counter-Strike 2 makes it possible to distinguish legitimate players from suspicious profiles using labeled datasets. However, these approaches also present limitations: the need for large volumes of data, the risk of false positives, delayed detection, and the continuous adaptation of cheaters. The objective of this thesis is therefore twofold. First, to analyze the advantages and limitations of kernel-level anti-cheats in light of cybersecurity, system stability, and user trust considerations. Second, to design and evaluate a behavioral detection system based on gameplay metrics. A proof of concept will be implemented using replay data or simulated logs for a Counter-Strike-like FPS in order to generate a suspicion score and study the prioritization of human review.

Finally, the thesis will propose an architectural positioning of these approaches within a multi-layer anti-cheat system, combining server-side detection, lightweight client mechanisms, and human validation, with the aim of reconciling effectiveness, security, and player acceptability.

\section{Outline}

In the first chapter, we will present the state of the art in cheating and anti-cheat technologies. We will analyze the main categories of cheats - automation tools, information cheats, and advanced hardware-assisted vectors - and examine the anti-cheat mechanisms developed in response, from client-side solutions to kernel-level \glspl{driver} and data-driven detection methods.

In the second chapter, we will assess the risks and structural limitations of kernel-level anti-cheats. We will analyze the technological arms race they sustain, their inability to address hardware-assisted cheating such as \gls{DMA}, and their economic and operational costs. This analysis will be supported by an experimental section using Windows-Kernel-Explorer to illustrate the level of system access a privileged kernel component holds, and to discuss the security and privacy implications that follow.

In the third chapter, we will present the design and implementation of a server-side gameplay analytics pipeline for suspicion scoring. We will describe the telemetry data model, the algorithmic architecture combining rule-based heuristics and an \gls{LSTM} layer, and the results of benchmarking against simulated player profiles. We will evaluate the system's performance, discuss the high-skill bias problem and its mitigations, and compare this approach against the kernel-level paradigm examined in the previous chapter.

Finally, in the conclusion, we will synthesize our findings and propose a positioning of these complementary approaches within a multi-layered anti-cheat architecture, balancing detection effectiveness, system security, and player trust.